% Reusable examples. Copy the component you need into a section file.

% Citation and cross-reference:
% Prior work established the core approach \cite{example2026}.
% See Section~\ref{section:methods} for the evaluation protocol.

% Equation:
% \begin{equation}
%   y = f_{\theta}(x),
%   \label{eq:example}
% \end{equation}

% Figure:
% \begin{figure}[htbp]
%   \centering
%   \includegraphics[width=0.8\linewidth]{example-figure.pdf}
%   \caption{Descriptive figure caption.}
%   \label{fig:example}
% \end{figure}

% Full-width figure in two-column mode:
% \begin{figure*}[t]
%   \centering
%   \includegraphics[width=0.9\textwidth]{example-wide-figure.pdf}
%   \caption{Descriptive full-width figure caption.}
%   \label{fig:example-wide}
% \end{figure*}

% Table:
% \begin{table}[htbp]
%   \centering
%   \caption{Descriptive table caption.}
%   \label{tab:example}
%   \begin{tabular}{lcc}
%     \toprule
%     Method & Metric A & Metric B \\
%     \midrule
%     Baseline & 0.00 & 0.00 \\
%     Proposed & 0.00 & 0.00 \\
%     \bottomrule
%   \end{tabular}
% \end{table}

% Full-width table in two-column mode:
% \begin{table*}[t]
%   \centering
%   \caption{Descriptive full-width table caption.}
%   \label{tab:example-wide}
%   \begin{threeparttable}
%     \begin{tabularx}{\textwidth}{Yccc}
%       \toprule
%       Method & Metric A & Metric B & Metric C \\
%       \midrule
%       Baseline & 0.00 & 0.00 & 0.00 \\
%       Proposed & 0.00 & 0.00 & 0.00 \\
%       \bottomrule
%     \end{tabularx}
%     \begin{tablenotes}[flushleft]\footnotesize
%       \item Note: define abbreviations and statistical annotations here.
%     \end{tablenotes}
%   \end{threeparttable}
% \end{table*}

% Theorem and proof:
% \begin{theorem}
%   State the result precisely.
% \end{theorem}
% \begin{proof}
%   Give the proof.
% \end{proof}

% Algorithm:
% \begin{algorithm}[htbp]
%   \caption{Example procedure}
%   \label{alg:example}
%   \begin{algorithmic}[1]
%     \Require Input data $X$
%     \Ensure Output $Y$
%     \State Initialize parameters
%     \State \Return $Y$
%   \end{algorithmic}
% \end{algorithm}

% Code listing:
% \begin{lstlisting}[language=Python,caption={Example code},label={lst:example}]
% def example(value):
%     return value
% \end{lstlisting}
